\documentclass[10pt,a4paper]{amsart}
\usepackage{amsmath,amssymb,amsthm,graphicx}
% Retained minimal preamble for the Stacks Project chapter sites-cohomology.tex (Cohomology on Sites).
% Verbatim excerpt of the chapter preamble: Xy-pic setup (native preamble lines 14-18) and the native
% theorem declarations with the mathematical macros (native preamble lines 154-217). Copied mechanically;
% nothing added, removed or redefined. The wrapper class and the amsmath/amssymb/amsthm/graphicx base
% packages are supplied by the assembler, so the class/externaldocument part of the native preamble is omitted.
% For commutative diagrams we use Xy-pic
\usepackage[all]{xy}

% We use 2cell for 2-commutative diagrams.
\xyoption{2cell}
% Theorem environments.
%
\theoremstyle{plain}
\newtheorem{theorem}[subsection]{Theorem}
\newtheorem{proposition}[subsection]{Proposition}
\newtheorem{lemma}[subsection]{Lemma}

\theoremstyle{definition}
\newtheorem{definition}[subsection]{Definition}
\newtheorem{example}[subsection]{Example}
\newtheorem{exercise}[subsection]{Exercise}
\newtheorem{situation}[subsection]{Situation}

\theoremstyle{remark}
\newtheorem{remark}[subsection]{Remark}
\newtheorem{remarks}[subsection]{Remarks}

\numberwithin{equation}{subsection}

% Macros
%
\def\lim{\mathop{\mathrm{lim}}\nolimits}
\def\colim{\mathop{\mathrm{colim}}\nolimits}
\def\Spec{\mathop{\mathrm{Spec}}}
\def\Hom{\mathop{\mathrm{Hom}}\nolimits}
\def\Ext{\mathop{\mathrm{Ext}}\nolimits}
\def\SheafHom{\mathop{\mathcal{H}\!\mathit{om}}\nolimits}
\def\SheafExt{\mathop{\mathcal{E}\!\mathit{xt}}\nolimits}
\def\Sch{\mathit{Sch}}
\def\Mor{\mathop{\mathrm{Mor}}\nolimits}
\def\Ob{\mathop{\mathrm{Ob}}\nolimits}
\def\Sh{\mathop{\mathit{Sh}}\nolimits}
\def\NL{\mathop{N\!L}\nolimits}
\def\CH{\mathop{\mathrm{CH}}\nolimits}
\def\proetale{{pro\text{-}\acute{e}tale}}
\def\etale{{\acute{e}tale}}
\def\QCoh{\mathit{QCoh}}
\def\Ker{\mathop{\mathrm{Ker}}}
\def\Im{\mathop{\mathrm{Im}}}
\def\Coker{\mathop{\mathrm{Coker}}}
\def\Coim{\mathop{\mathrm{Coim}}}

% Boxtimes
%
\DeclareMathSymbol{\boxtimes}{\mathbin}{AMSa}{"02}

%
% Macros for moduli stacks/spaces
%
\def\QCohstack{\mathcal{QC}\!\mathit{oh}}
\def\Cohstack{\mathcal{C}\!\mathit{oh}}
\def\Spacesstack{\mathcal{S}\!\mathit{paces}}
\def\Quotfunctor{\mathrm{Quot}}
\def\Hilbfunctor{\mathrm{Hilb}}
\def\Curvesstack{\mathcal{C}\!\mathit{urves}}
\def\Polarizedstack{\mathcal{P}\!\mathit{olarized}}
\def\Complexesstack{\mathcal{C}\!\mathit{omplexes}}
% \Pic is the operator that assigns to X its picard group, usage \Pic(X)
% \Picardstack_{X/B} denotes the Picard stack of X over B
% \Picardfunctor_{X/B} denotes the Picard functor of X over B
\def\Pic{\mathop{\mathrm{Pic}}\nolimits}
\def\Picardstack{\mathcal{P}\!\mathit{ic}}
\def\Picardfunctor{\mathrm{Pic}}
\def\Deformationcategory{\mathcal{D}\!\mathit{ef}}
\title{Boundedness above and tor amplitude survive passage from O/I to O/I\textasciicircum{}n}
\author{The Stacks Project Authors}
\date{stacks-project:sites-cohomology.tex:local-snapshot-20261011; GNU FDL 1.2 or later}
\begin{document}
\maketitle
\noindent\textit{Editorial note.} All delivered mathematical text is the author's own, copied from the pinned Stacks Project chapter file sites-cohomology.tex; no statement or proof step was rewritten, reordered or removed.\par\medskip
\noindent\textit{Editorial note.} The proof is the author's complete argument for both parts. The induction of part (1) is written as an induction without spelling out the base case, and the counted span keeps that wording.\par\medskip
\noindent\textit{Editorial note.} The label lemma-bounded occurs in no other staged Stacks chapter, so no cross-chapter restatement had to be weighed for this card.\par\medskip
\noindent\textit{Editorial note.} The preamble is the chapter's own retained minimal preamble (sources/stacks-sites-cohomology/retained-preamble.tex), a verbatim excerpt of the native preamble: the Xy-pic setup and the native theorem declarations with the mathematical macros.\par\medskip
\noindent\textit{Allowed standard prerequisites (editorial).}\begin{itemize}
\item Derived category of modules over a ringed site and the derived tensor product. Same chapter, the unbounded derived category D(O) of a ringed site and the derived tensor product of complexes of O-modules. Supplies the objects K tensor\textasciicircum{}L O/I\textasciicircum{}n and their cohomology sheaves, which are the quantities bounded in both parts.
\item Tor amplitude of an object of D(O/I). Same chapter, the definition of tor amplitude in [a, b]: the vanishing of H\textasciicircum{}i(K tensor\textasciicircum{}L G) for i outside [a, b] and every O/I-module G. Supplies the property transferred in part (2), and the reduction of that part to the displayed vanishing of cohomology sheaves.
\item Distinguished triangles from the short exact sequence of ideals I\textasciicircum{}n/I\textasciicircum{}\{n+1\}. Same chapter, the distinguished triangles attached to the short exact sequence 0 -> I\textasciicircum{}n/I\textasciicircum{}\{n+1\} -> O/I\textasciicircum{}\{n+1\} -> O/I\textasciicircum{}n -> 0 after applying K tensor\textasciicircum{}L -. Supplies the inductive triangles used in part (1) to control the cohomology of K tensor\textasciicircum{}L O/I\textasciicircum{}n uniformly in n.
\item Finite filtration of a module over O/I\textasciicircum{}n by modules over O/I. Same chapter, the filtration 0 subset I\textasciicircum{}\{n-1\}F subset ... subset IF subset F whose successive quotients are O/I-modules. Supplies the reduction in part (2) from arbitrary O/I\textasciicircum{}n-modules to O/I-modules.
\end{itemize}\medskip
\begin{lemma}
\label{lemma-bounded}
Let $(\mathcal{C}, \mathcal{O})$ be a ringed site.
Let $\mathcal{I} \subset \mathcal{O}$ be a sheaf of ideals.
Let $K$ be an object of $D(\mathcal{O})$.
\begin{enumerate}
\item If $K \otimes_\mathcal{O}^\mathbf{L} \mathcal{O}/\mathcal{I}$
is bounded above, then
$K \otimes_\mathcal{O}^\mathbf{L} \mathcal{O}/\mathcal{I}^n$
is uniformly bounded above for all $n$.
\item If $K \otimes_\mathcal{O}^\mathbf{L} \mathcal{O}/\mathcal{I}$
as an object of $D(\mathcal{O}/\mathcal{I})$ has tor amplitude in $[a, b]$,
then $K \otimes_\mathcal{O}^\mathbf{L} \mathcal{O}/\mathcal{I}^n$
as an object of $D(\mathcal{O}/\mathcal{I}^n)$
has tor amplitude in $[a, b]$ for all $n$.
\end{enumerate}
\end{lemma}

\begin{proof}
Proof of (1). Assume that
$K \otimes_\mathcal{O}^\mathbf{L} \mathcal{O}/\mathcal{I}$
is bounded above, say
$H^i(K \otimes_\mathcal{O}^\mathbf{L} \mathcal{O}/\mathcal{I}) = 0$
for $i > b$. Note that we have distinguished triangles
$$
K \otimes_\mathcal{O}^\mathbf{L}
\mathcal{I}^n/\mathcal{I}^{n + 1} \to
K \otimes_\mathcal{O}^\mathbf{L}
\mathcal{O}/\mathcal{I}^{n + 1} \to
K \otimes_\mathcal{O}^\mathbf{L}
\mathcal{O}/\mathcal{I}^n \to
K \otimes_\mathcal{O}^\mathbf{L}
\mathcal{I}^n/\mathcal{I}^{n + 1}[1]
$$
and that
$$
K \otimes_\mathcal{O}^\mathbf{L}
\mathcal{I}^n/\mathcal{I}^{n + 1} =
\left(
K \otimes_\mathcal{O}^\mathbf{L}
\mathcal{O}/\mathcal{I}\right)
\otimes_{\mathcal{O}/\mathcal{I}}^\mathbf{L}
\mathcal{I}^n/\mathcal{I}^{n + 1}
$$
By induction we conclude that
$H^i(K \otimes_\mathcal{O}^\mathbf{L} \mathcal{O}/\mathcal{I}^n) = 0$
for $i > b$ for all $n$.

\medskip\noindent
Proof of (2). Assume $K \otimes_\mathcal{O}^\mathbf{L} \mathcal{O}/\mathcal{I}$
as an object of $D(\mathcal{O}/\mathcal{I})$ has tor amplitude in $[a, b]$.
Let $\mathcal{F}$ be a sheaf of $\mathcal{O}/\mathcal{I}^n$-modules.
Then we have a finite filtration
$$
0 \subset \mathcal{I}^{n - 1}\mathcal{F} \subset \ldots
\subset \mathcal{I}\mathcal{F} \subset \mathcal{F}
$$
whose successive quotients are sheaves of $\mathcal{O}/\mathcal{I}$-modules.
Thus to prove that $K \otimes_\mathcal{O}^\mathbf{L} \mathcal{O}/\mathcal{I}^n$
has tor amplitude in $[a, b]$ it suffices to show
$H^i(K \otimes_\mathcal{O}^\mathbf{L} \mathcal{O}/\mathcal{I}^n
\otimes_{\mathcal{O}/\mathcal{I}^n}^\mathbf{L} \mathcal{G})$
is zero for $i \not \in [a, b]$ for all $\mathcal{O}/\mathcal{I}$-modules
$\mathcal{G}$. Since
$$
\left(K \otimes_\mathcal{O}^\mathbf{L} \mathcal{O}/\mathcal{I}^n\right)
\otimes_{\mathcal{O}/\mathcal{I}^n}^\mathbf{L} \mathcal{G}
=
\left(K \otimes_\mathcal{O}^\mathbf{L} \mathcal{O}/\mathcal{I}\right)
\otimes_{\mathcal{O}/\mathcal{I}}^\mathbf{L} \mathcal{G}
$$
for every sheaf of $\mathcal{O}/\mathcal{I}$-modules $\mathcal{G}$
the result follows.
\end{proof}

\end{document}
